% 第7章 初期条件と原始ゆらぎ（WP5 skeleton; 本文・導出は執筆予定）
\section{初期条件と原始ゆらぎ}
\label{ch:ch07_initial}

\paragraph{この章で導くこと（入力→出力）}
超地平線極限で断熱初期条件一式を導き、$\mathcal R$ 保存と規格化、原始スペクトルを確立する。

% --- 節構成（01_textbook_spec.md より）-------------------------------------
\subsection{超地平線極限での縮約}
% TODO(WP5): 7.1 の本文。
\subsection{断熱条件と全初期値（契約 S7.1）}
% TODO(WP5): 7.2 の本文。
\subsection{$\mathcal R$ の保存と規格化（契約 S7.2）}
% TODO(WP5): 7.3 の本文。
\subsection{インフレーションの最小限}
% TODO(WP5): 7.4 の本文。
\subsection{ガウシアン確率場とパワースペクトル}
% TODO(WP5): 7.5 の本文。


\paragraph{導出契約}
本章の導出契約: S7.1–S7.2。各契約は「入力式→出力式」を中間ステップ省略なし
（1ステップ＝学部4年生が5分で検算可能）で履行する。\emph{本文プローズは WP5 で執筆。}

\paragraph{まとめ・チェックリスト・演習}
% TODO(WP5): 章末まとめ、チェックリスト、演習 3–6 問（うち1問は対応 NB の課題）。
